\documentclass[conference]{IEEEtran}
\IEEEoverridecommandlockouts
\usepackage{cite}
\usepackage{amsmath,amssymb,amsfonts}
\usepackage{algorithmic}
\usepackage{graphicx}
\usepackage{textcomp}
\usepackage{xcolor}
\usepackage{url}
\usepackage{hyperref}

            
\def\BibTeX{{\rm B\kern-.05em{\sc i\kern-.025em b}\kern-.08em
    T\kern-.1667em\lower.7ex\hbox{E}\kern-.125emX}}
\begin{document}


%===========================================
% Start your report from here.
%===========================================


\title{Your Title\\
{\footnotesize Data Mining Final Project, 2026 Spring}
}

\author{\IEEEauthorblockN{Member 1's name}
\and
\IEEEauthorblockN{Member 2's name}
\and
\IEEEauthorblockN{Member 3's name}
}

\maketitle

\textcolor{blue}{This template is built on IEEE paper template\footnote{\url{https://www.overleaf.com/latex/templates/ieee-conference-template/grfzhhncsfqn}}}


\begin{abstract}
Please provide your group information (Group ID) and your GitHub link here. For example: Our code is available at GitHub:~\href{https://github.com/NYCU-Data-Mining/DM2026-Assignment-1}{\color{blue}https://github.com/NYCU-Data-Mining/DM2026-Assignment-1}.


\end{abstract}


\section{Project Summary (\textcolor{red}{max 1 page)}}

\begin{itemize}
    \item Describe the goal of this project
    \item Provide the high-level idea of your method
    \item Introduce the related works here
\end{itemize}


Here are some papers~\cite{iglovikov2017satellite}\cite{kechyn2018sales} related to other Kaggle competitions for your reference.


\section{Proposed Method (\textcolor{red}{1-3 pages)}}

Please describe your proposed method in detail. You may include
\begin{itemize}
    \item Formulas
    \item Pseudocode
    \item Framework diagrams
    \item or other materials to enrich the content.
\end{itemize}


\section{Experiments (\textcolor{red}{2-4 pages)}}

You may consider structuring this section as follows:

\subsubsection{Dataset and Experimental Setup}

You may include the data description, statistics of data, and the basic setup for your experiments.

\subsubsection{Baseline Descriptions}

You can include your self-defined baselines to reveal the robustness of your method.


\subsubsection{Experimental Results}

We encourage you to include experiments and analysis from different perspectives, such as:

\begin{itemize}
    \item Your model's performance on Kaggle
    \item Comparison with your self-defined baselines
    \item Ablation study, Parameter analysis, Case study
    \item ...
\end{itemize}

You are free to include any figures and tables to present your experiments.


%===========================================
% Reference
%===========================================

\bibliographystyle{IEEEtran}  
\bibliography{reference}  

\end{document}
